\subsubsection{Kuhn (bipartite matching)}
\begin{lstlisting}[style=mystyle, language=C++]
// TC: O(V * E) donde V = nodos lado izquierdo, E = aristas
// usa: matching maximo en grafo bipartito; usa DFS con vis para encontrar augmenting paths
// nota: para N, M <= 1e3 anda bien; para mas usar Hopcroft-Karp (O(sqrt(V) * E))
#include <bits/stdc++.h>
using namespace std;

int n, m;
vector<int> adj[1005];

bool try_match(int node, vector<int> adj[], vector<int>& matchR, vector<bool>& vis){
	for(auto adjN : adj[node]){
		if(vis[adjN]) continue;
		vis[adjN] = true;
		if(matchR[adjN] == -1 || try_match(matchR[adjN], adj, matchR, vis)){
			matchR[adjN] = node;
			return true;
		}
	}
	return false;
}

int main(){
	// entrada: n nodos izquierda, m nodos derecha, m aristas (u, v) con u en [1..n], v en [1..m]
	cin >> n >> m;
	for(int i=0 ; i<m ; i++){
		int u, v; cin >> u >> v;
		adj[u].push_back(v);
	}
	vector<int> matchR(m + 1, -1);
	int matching = 0;
	for(int node = 1; node <= n; node++){
		vector<bool> vis(m + 1, false);
		if(try_match(node, adj, matchR, vis)) matching++;
	}
	cout << matching << "\n";
	return 0;
}
\end{lstlisting}
